\uuid{8PaF}
\exo7id{6385}
\titre{exo7 6385}
\auteur{potyag}
\organisation{exo7}
\datecreate{2011-10-16}
\isIndication{false}
\isCorrection{false}
\chapitre{Isométrie euclidienne}
\sousChapitre{Isométrie euclidienne}
\cpge{mpsi-mp2i.calcul-matriciel-systemes.a, pcsi.calcul-matriciel-systemes-lineaires.a, ptsi.calcul-matriciel-systemes-lineaires.a, mpsi-mp2i.determinants.d, pcsi.calcul-matriciel-systemes-lineaires.d, ptsi.calcul-matriciel-systemes-lineaires.d, mp-mpi.endomorphismes-euclidien.b, pc.endomorphismes-euclidiens.b, psi.endomorphismes-euclidien.b, pt.isometries-euclidien.b, mp-mpi.endomorphismes-euclidien.c, pc.endomorphismes-euclidiens.a, psi.endomorphismes-euclidien.a, pt.isometries-euclidien.a}
\module{Algèbre et géométrie}
\niveau{L3}
\difficulte{}

\contenu{
\texte{
Soit $E$ un espace vectoriel euclidien de dimension $n$.
}
\begin{enumerate}
    \item \question{Montrer que $A\in GL(E)$ appartient à $\mathcal{O}(n)$ si et seulement si
${ }^T AA =I$.}
    \item \question{Montrer que si $A\in \mathcal{O}(n)$ alors $\det A = \pm 1$.}
    \item \question{Montrer que $A \in U(n)$ si et seulement si ${}^T A\bar A =I$.}
\end{enumerate}
}
